\section{Source action, exterior boundary response and optical normalization}
\label{app:exterior-derivation}

This section defines the precise conditional exterior problem used by the
tracker. The action, background and source pairing are supplied; a physical
black-hole horizon relation and the microscopic BHSM-to-R1 identification
are not inferred. The worldtube below is outside the unresolved compact
interior and is not a horizon boundary condition.

\subsection{Common action and local source}

In signature $(-+++)$, initially setting $c=\hbar=1$ and
$M_{\rm Pl}^{-2}=8\pi G$, the retained action is
\begin{equation}
 S=\int\sqrt{-g}\left[\frac{M_{\rm Pl}^2}{2}R
 -\frac12(\nabla T)^2+\frac{(\nabla T)^2\Box T}{\lambda}
 -V_0+\mu^3T\right]d^4x+S_m .
 \label{eq:ext-action}
\end{equation}
Here physical $\lambda=\Lambda^3$, while the pinned numerical input is
$\widehat\lambda=\lambda/(M_{\rm Pl}H_0^2)=1$. This is a cubic
kinetic-braiding action in the established scalar--tensor setting
\cite{Deffayet2010,Kimura2012}. With
$\Pi_{\mu\nu}=\nabla_\mu\nabla_\nu T$, variation gives
\begin{equation}
 \Box T+\frac2\lambda\left[\Pi_{\mu\nu}\Pi^{\mu\nu}
 -(\Box T)^2+R_{\mu\nu}\nabla^\mu T\nabla^\nu T\right]+\mu^3=0.
 \label{eq:ext-scalar-euler}
\end{equation}
It reproduces the frozen homogeneous current equation
$J=v+6Hv^2/\lambda$, $\dot J+3HJ=\mu^3$, where $v=\dot{\bar T}$.
Thus the local source sign follows from the same action and background.

For a minimally coupled point-matter model,
$S_{\rm pp}=-M\int ds$ and
$\delta S_{\rm pp}=\tfrac12\int\sqrt{-g}T^{\mu\nu}_{\rm pp}
\delta g_{\mu\nu}d^4x$. In its nonrelativistic local limit the excess
density is $M\delta^{(3)}(\boldsymbol r)$, with no direct scalar insertion
$\delta S_{\rm pp}/\delta T$. Strongly self-gravitating horizon response
does not follow from treating that point-matter action as a black hole.

Subtract the homogeneous equations, write $T=\bar T(t)+\pi_T(r)$ and
freeze $v$, $d=\ddot{\bar T}$ and $H$ during a local encounter. At the
retained weak-metric, quasi-static order, keeping the nonlinear scalar
Hessian terms, the scalar and metric rows are
\begin{align}
 A\Delta\pi_T-\frac2\lambda
 [(\Delta\pi_T)^2-|\operatorname{Hess}\pi_T|^2]+b\Delta\Phi&=0,
 \nonumber\\
 2M_{\rm Pl}^2\Delta\Phi&=\delta\rho-b\Delta\pi_T,
 \qquad\Phi=\Psi,\nonumber\\
 b=\frac{2v^2}{\lambda},\qquad
 A&=1+\frac{4(d+2Hv)}{\lambda}.
 \label{eq:ext-qs-rows}
\end{align}
The density is the local excess, not the homogeneous cosmic density
counted again. Eliminating that same metric row gives
\begin{equation}
 C\Delta\pi_T-\frac2\lambda
 [(\Delta\pi_T)^2-|\operatorname{Hess}\pi_T|^2]
 =-\frac{b\,\delta\rho}{2M_{\rm Pl}^2},\qquad
 C=A-\frac{b^2}{2M_{\rm Pl}^2}.
 \label{eq:ext-debraided}
\end{equation}
For example, varying the static functional
\begin{equation}
 S_{\rm QS}=\int dt\,d^3r\left[-\frac C2|\nabla\pi_T|^2
 +\frac{|\nabla\pi_T|^2\Delta\pi_T}{\lambda}
 +\frac{b\,\delta\rho\,\pi_T}{2M_{\rm Pl}^2}\right]
 \label{eq:ext-qs-functional}
\end{equation}
reproduces Eq.~\eqref{eq:ext-debraided}. Its last term is the consumed
source pairing after metric elimination, not an additional covariant
matter coupling. Higher weak-field orders need not retain $\Phi=\Psi$.

\subsection{Spherical exterior, scales and open compact-object flux}

Let $L_0=c/H_0$, $x=r/L_0$, $\widehat m=(GM/c^2)/L_0$,
$\widehat v=v/(M_{\rm Pl}H_0)$, and
$\widehat d=d/(M_{\rm Pl}H_0^2)$. Define
\begin{equation}
 \widehat\beta=\frac b{M_{\rm Pl}}
 =\frac{2\widehat v^2}{\widehat\lambda},\quad
 C=1+\frac{4(\widehat d+2\widehat H\widehat v)}{\widehat\lambda}
 -\frac{\widehat\beta^2}{2},\qquad
 p(x)=\frac{d(\pi_T/M_{\rm Pl})}{dx}.
 \label{eq:ext-beta}
\end{equation}
Regular-origin matching for an ordinary finite matter source gives the
exterior flux $-\widehat\beta\widehat m$. To expose the uncompleted
compact-object matching, retain a supplied nonnegative flux ratio $\chi$:
\begin{align}
 Cx^2p-\frac4{\widehat\lambda}xp^2
 &=-\chi\widehat\beta\widehat m,\nonumber\\
 p(x)&=\frac{-2\chi\widehat\beta\widehat m}
 {x^2\left[C+\sqrt{C^2+16\chi\widehat\beta\widehat m/
 (\widehat\lambda x^3)}\right]},\nonumber\\
 \Phi_x=\Psi_x&=\frac{\widehat m}{x^2}
 -\frac{\widehat\beta}{2}p(x).
 \label{eq:ext-radial-branch}
\end{align}
This root is continuous with the decaying linear response. The metric
monopole mass remains $M$; $\chi$ changes the scalar flux, not the mass.
The value $\chi=1$ is the regular-matter reference, not a derived
black-hole charge. A horizon-regular, cosmologically matched solution is
required to determine the physical horizon relation. A negative flux is
outside the real, positive-discriminant branch implemented here.

The exterior screening scale and extra force fraction are
\begin{equation}
 r_V^3=\frac{16\chi\widehat\beta\,(GM/c^2)L_0^2}
 {\widehat\lambda C^2},\quad
 \epsilon_\infty=\frac{\chi\widehat\beta^2}{2C},\quad
 F(r)=\frac{\delta\Phi'}{(GM/c^2)/r^2}
 =\frac{2\epsilon_\infty}{1+\sqrt{1+(r_V/r)^3}}.
 \label{eq:ext-screening}
\end{equation}
The scalar field is not the provider's time-displacement column:
$\pi_T=-v\pi_{\rm time}$ when reconstructing it from uniform-field gauge.
Adding it directly as another gravitational potential would change both
its normalization and the source equation.

At $z_l=0.3$ the pinned background replay gives
\begin{equation}
 \widehat\beta=0.113063844895134,
 \quad C=4.23885855943370,
 \quad D_t=4.39622102870814,
 \label{eq:ext-coeff-values}
\end{equation}
where $D_t=1+12\widehat H\widehat v/\widehat\lambda
+6\widehat v^4/\widehat\lambda^2$ is the homogeneous leading time
principal coefficient. For $M=10^9M_\odot$, $H_0=67.4$ km/s/Mpc,
$\chi=1$ and $r_0=10$ pc,
$(L_0,GM/c^2,r_V)=(4.4479593175\times10^9,
4.7854158416\times10^{-5},45680.3117591)$ pc.
The compactness at the worldtube is $4.7854\times10^{-6}$ and
$H_lr_V/c=1.2156\times10^{-5}$. These are local expansion parameters,
not measured error bars. Near a horizon the weak-metric approximation
fails; the worldtube must remain outside that domain.

\subsection{Exact linear multipoles and canonical boundary response}

Write the dimensionless scalar perturbation as
$\delta\widehat\pi=a f_\ell(r)Y_{\ell m}$, with
$\int|Y_{\ell m}|^2d\Omega=1$. Linearizing
Eq.~\eqref{eq:ext-debraided} gives
\begin{equation}
 \delta E=Z^{ij}\partial_i\partial_j\delta\pi
 =\partial_i(Z^{ij}\partial_j\delta\pi),\qquad
 Z^{ij}=C\delta^{ij}-\frac4\lambda
 [(\Delta\pi_0)\delta^{ij}-\Pi_0^{ij}],
 \label{eq:ext-linear-operator}
\end{equation}
because $\partial_iZ^{ij}=0$. The anisotropic Hessian term is essential.
The spherical eigenvalues satisfy
\begin{equation}
 K_r=C-\frac{8\pi_0'}{\lambda r},\quad
 K_\Omega=C-\frac4\lambda(\pi_0''+\pi_0'/r),\quad
 K_r'=\frac{2(K_\Omega-K_r)}r.
\end{equation}
Outside the source the normalized radial equation is therefore
\begin{equation}
 \frac1{r^2}(r^2K_rf_\ell')'
 -\frac{\ell(\ell+1)K_\Omega}{r^2}f_\ell=0.
 \label{eq:ext-multipole}
\end{equation}
An actual density multipole adds the fixed source
$-b\delta\rho_{\ell m}/(2M_{\rm Pl}^2)$ to the right side and requires
interior matching. A prescribed boundary value is not that interior solve.

For $\rho=r/r_V$ and $u=\rho^{-3}$, the exact exterior coefficients are
\begin{equation}
 K_r=C\sqrt{1+u},\quad
 K_\Omega=C\frac{1+u/4}{\sqrt{1+u}},\quad
 K_t=D_t+\frac{3(S-C)^2}{4S}\ge D_t>0,
 \quad S=C\sqrt{1+u}.
 \label{eq:ext-principal}
\end{equation}
Thus this retained scalar problem is elliptic and has a positive leading
time coefficient. This is not full-system stability or a horizon/UV
certificate; its scalar characteristic speeds can exceed the metric
light speed in part of the exterior.

The decaying solution of Eq.~\eqref{eq:ext-multipole} is
\begin{equation}
 R_\ell(\rho)=\rho^{-\ell-1}
 {}_2F_1\!\left(\frac{\ell+1}{2},\frac{\ell+2}{6};
 \frac{2\ell+4}{3};-\rho^{-3}\right),\qquad
 f_\ell(r)=\frac{R_\ell(r/r_V)}{R_\ell(r_0/r_V)}.
 \label{eq:ext-hypergeometric}
\end{equation}
It follows by substituting $t=-\rho^{-3}$ in the radial equation.
The deep-screened powers are $(\ell+1)/2$ and $-\ell/2$;
the far powers are $\ell$ and $-\ell-1$. In particular the decaying
quadrupole has $r^{-1}$ and $r^{-3}$ limits. For $\ell=1$,
${}_2F_1(1,1/2;2;-u)=2/(1+\sqrt{1+u})$, so $f_1$ is exactly
proportional to $-p$: translating the monopole reproduces the linear
dipole equation. This supplies an independent field-equation check.
The $\chi=0$ limit is $f_\ell=(r_0/r)^{\ell+1}$.

The exterior quadratic energy, divided by its overall
$M_{\rm Pl}^2L_0$ factor, is
\begin{equation}
 \frac12\int_{x_0}^\infty
 [x^2K_r(\partial_x\delta\widehat\pi)^2
 +\ell(\ell+1)K_\Omega\delta\widehat\pi^2]dx
 =\frac12\Lambda_\ell a^2,\qquad
 \Lambda_\ell=-x_0^2K_r(x_0)f_{\ell,x}(x_0)>0.
 \label{eq:ext-impedance}
\end{equation}
The inner normal of the exterior domain is $-\partial_r$; its conjugate
flux is $q_{\ell m}=\Lambda_\ell a$. Positivity establishes uniqueness
of the zero-data decaying exterior solution. Holding $a$ fixed and
holding $q_{\ell m}$ fixed are different source comparisons. Neither
amplitude is selected by the action without an interior/boundary state.
The numerical reference has $\Lambda_2=0.00294226449484$ in these
dimensionless-$x$ conventions.

The same metric constraint yields
$\delta\Phi=\delta\Psi=\delta\Phi_N-
\widehat\beta\delta\widehat\pi/2+H_{\rm match}$, where
$\Delta H_{\rm match}=0$. Ordinary gravitational multipoles and outer
matching remain independent until a common source fixes them. The
implemented scalar-driven term adds no arbitrary metric quadrupole.

\subsection{Projected phase, derivatives and finite matching}

For a transverse real quadrupole choose
\begin{equation}
 Y_N=\sqrt{\frac{15}{16\pi}}(n_x^2-n_y^2)
 \equiv N_2(n_x^2-n_y^2),\qquad
 W\equiv\Phi+\Psi=-\widehat\beta a f_2(r)Y_N.
\end{equation}
For impact $b\ge r_0$, every integration point lies outside the worldtube.
The scalar projected physical phase length is
\begin{align}
 \delta\widehat\psi(b,\varphi)&=a\,q(b)\cos2\varphi,\nonumber\\
 q(b)&=-\widehat\beta N_2\int_{-\infty}^{\infty}
 f_2(\sqrt{b^2+z^2})\frac{b^2}{b^2+z^2}\,dz\nonumber\\
 &=-2\widehat\beta N_2 b\int_0^{\pi/2}f_2(b\sec\theta)d\theta.
 \label{eq:ext-phase}
\end{align}
Both halves of the line are included. With $r=b\sec\theta$,
\begin{align}
 q'(b)&=-2\widehat\beta N_2\int_0^{\pi/2}(f_2+rf_2')d\theta,\nonumber\\
 q''(b)&=-\frac{2\widehat\beta N_2}{b}
 \int_0^{\pi/2}(r^2f_2''+2rf_2')d\theta.
 \label{eq:ext-phase-derivatives}
\end{align}
The provider returns $q,q',q''$ in pc, dimensionless units and
pc$^{-1}$, respectively. Exact hypergeometric connection formulae
preserve the cancellation of the deep $r^{-1}$ contribution in these
radial derivatives. Projection gives a deep phase plateau and a far
$b^{-2}$ phase, rather than identifying a local force ratio with
a projected photon deflection.

An optional axial boundary coordinate uses the independently normalized
$Y_{20}=N_0(3n_z^2-1)$, $N_0=\sqrt{5/(16\pi)}$. Its physical phase
per unit axial boundary amplitude is
\begin{equation}
 q_0(b)=-2\widehat\beta N_0b\int_0^{\pi/2}
 (2\sec^2\theta-3)f_2(b\sec\theta)d\theta.
 \label{eq:ext-axial-phase}
\end{equation}
Its first and second radial derivatives replace $f_2$ by $f_2+rf_2'$
and $r^2f_2''+2rf_2'$, respectively, with prefactors
$-2\widehat\beta N_0$ and $-2\widehat\beta N_0/b$.
This supplies a separately normalized radial nuisance when the axial
state is not fixed. It is distinct from varying the spherical scalar
flux. Its leading deep-screened phase is logarithmic and can be
mass-like. In the exact zero-flux harmonic limit its full-line
projection vanishes; the transverse phase is then exactly
$q=-4\widehat\beta N_2r_0^3/(3b^2)$. The implementation uses these
analytic zero-flux identities instead of an interpolation residual.

The scalar monopole deflection is supplied separately:
\begin{equation}
 \delta\widehat\alpha_0(b)
 =-\frac{\widehat\beta b}{L_0}\int_{-\infty}^{\infty}
 \frac{p(r/L_0)}r\,dz
 =\frac{4GM}{c^2b}\int_0^{\pi/2}F(b\sec\theta)\cos\theta\,d\theta.
 \label{eq:ext-monopole-deflection}
\end{equation}
The ordinary $4GM/(c^2b)$ is counted once by the lens model. A relative
spherical phase is obtained by integrating this deflection between
impacts; its formally divergent absolute line potential is unnecessary.

Infinity in these expressions is a matched thin-lens extension, not an
exact static spacetime. A local overlap cutoff obeys
$\max(b,r_V)\ll L\ll c/H_l,R_{\rm curvature}$. The scalar monopole
tail, divided by $4GM/(c^2b)$, is bounded at leading order by
$\epsilon_\infty b^2/(2L^2)$. If
$A_2=r_V^3/R_2(r_0/r_V)$, the Euler integral of the hypergeometric
function gives $f_2(r)\le A_2/r^3$ and hence
\begin{equation}
 |q_{|z|>L}|\le
 \frac{\widehat\beta N_2 A_2 b^2}{2L^4}.
 \label{eq:ext-phase-tail}
\end{equation}
These control the stated tail only. They do not bound coefficient time
dependence, a changed environment, or a different physical outer state.

\subsection{Perturbative example and actual numerical control}

For the supplied state convention
$a=\eta x_0|p(x_0)|/6$, $\eta=10^{-3}$, the numerical reference has
$a=1.37809214552\times10^{-16}$. The closed R1 distances at
$(z_l,z_s)=(0.3,1)$ give an Einstein impact $339.621287478$ pc and
$a q(r_E)/(4GM/c^2)=-8.891115935\times10^{-13}$.
These amplitudes are declared inputs, not horizon-selected values.

On 65 radii from $r_0$ to $10r_V$ and 100 directions, the sampled
maximum perturbation/background scalar-Hessian norm ratio is
$3.742\times10^{-4}$. The exact linearized equation residual, scaled
by its principal Hessian terms, is below $6.1\times10^{-15}$. The positive
comparison scale used here is
\begin{equation}
 \mathcal S=\left[C+\frac4\lambda
 (|\operatorname{tr}\Pi_0|+\|\Pi_0\|_F)\right]\|\delta\Pi\|_F .
\end{equation}
The retained nonlinear scalar remainder is exactly
$-2[(\operatorname{tr}\delta\Pi)^2-\delta\Pi:\delta\Pi]/\lambda$;
its sampled ratio to the stated principal-term scale is at most
$8.57\times10^{-5}$. These are finite-sample diagnostics, not uniform
interval certificates.

Independent radial integration agrees with the special function to
$2.25\times10^{-10}$ relative over the tested domain, with the outer
decaying condition imposed at $1000r_V$. The exact translation check is
within $1.23\times10^{-15}$. Independent original-line quadrature
checks the projected quadrupole phase to $2.23\times10^{-16}$ pc and
the scalar monopole deflection to $4.04\times10^{-27}$ at $b=1000$ pc.
The code records Gauss-order refinement and Hermite interpolation
discrepancies for every prepared table. Its fast potential, gradient
and Hessian derive from one interpolating potential, preserving their
mutual consistency.

Slow source motion requires small velocity relative to the scalar
characteristic speeds; a varying multipole additionally requires a
quasi-static crossing condition such as $\omega r_V/c_{\rm scalar}\ll1$.
The retained field equations omit higher weak-metric and time-dependent
terms. A nonzero environmental scalar Hessian modifies the same
operator and cannot generally be added by linear superposition. The
paper's homogeneous-background exterior avoids that extra assumption.
The positive principal coefficients, small illustrative amplitude and
tiny formal delay do not establish physical precision at that delay,
an observational forecast, or a solved black-hole boundary problem.
